%% ch08.tex --- Neural Networks: Representation and Forward Propagation
%% Numbers are computed by verify/ch08_forward.py.
%% The 2-3-1 network (N231) is the same one used in the companion lecture
%% slides, so that every number can be checked against either source.

\chapter{Neural Networks: Representation and Forward Propagation}
\label{ch:ch08}
\chaptermark{Neural Networks: Forward Propagation}

\epigraph{The question of whether a computer can think is no more interesting
than the question of whether a submarine can swim.}{---\ Edsger W.~Dijkstra}

\begin{chapterlearning}
After working through this chapter you will be able to:
\begin{enumerate}
  \item Describe the artificial neuron, say what it borrows from the biological
        one, and say where the analogy stops being useful.
  \item Explain the difference between engineering features by hand and
        learning them, and say what the designer of a network still chooses.
  \item Write the computation of one layer in vector and matrix form, using the
        layer notation $a^{[l]}_j$, and count its parameters.
  \item Carry out forward propagation through a network by hand, then as a
        sequence of matrix products over a whole batch.
  \item Show that one logistic unit cannot compute XOR and that two layers can,
        and explain what the hidden layer did to make the problem easy.
  \item Prove that a network with linear activations collapses to a single
        linear layer.
  \item State the universal approximation theorem and what it does and does not
        promise.
  \item Build the same network in Keras and in NumPy, with the shape conventions
        that make the two agree.
\end{enumerate}
\end{chapterlearning}

Every model so far has been linear in its parameters, and every non-linear
decision boundary has come from features we built by hand: $x^2$ in
\cref{sec:polynomial}, $x_1^2 + x_2^2$ in \cref{sec:decision-boundary}, all
twenty-seven monomials in \cref{sec:reg-logistic}. That works when we know which
features to build. It fails when the input is a photograph with a million
pixels, none of which means anything on its own, and nobody can say which
combinations of them matter.

A neural network learns its own features. It does so by stacking the logistic
unit of \cref{ch:ch06}: the first layer computes features of the input, the next
computes features of those features, and the last is a logistic regression on
whatever the layers below have found. This chapter defines the model and shows
how it computes a prediction --- \emph{forward propagation}. \Cref{ch:ch09} shows
how its weights are learned.

\section{Neurons and the Brain}
\label{sec:neurons-brain}

The original motivation for neural networks was to build software that learned
the way brains do. A biological neuron receives electrical signals from other
neurons through its \emph{dendrites}, integrates them in its cell body, and, if
the combined input is strong enough, fires a signal down its \emph{axon} to the
neurons it connects to. The artificial neuron keeps the outline of that story
and discards almost everything else.

\begin{defbox}{Artificial neuron (unit)}
A \term{unit} receives a vector of inputs $\va$, forms a weighted sum plus a
bias, and passes the result through an \term{activation function} $\gsig$:
\[
  z = \vw\cdot\va + b, \qquad a_{\text{out}} = \gsig(z) .
\]
The output $a_{\text{out}}$ is the unit's \term{activation} --- the number it
sends on. With $\gsig$ the sigmoid, a single unit is exactly the logistic
regression model of \cref{ch:ch06}.
\end{defbox}

\begin{figure}[htbp]
  \centering
  \begin{tikzpicture}[
      nninp/.style={circle, draw=MLTeal, fill=MLTeal!12, thick, minimum size=7mm,
                  inner sep=0pt, font=\small},
      nnunit/.style={circle, draw=MLBlue, fill=MLBlue!8, thick, minimum size=13mm,
                   inner sep=0pt, font=\small},
      nnsyn/.style={-{Stealth[length=2mm]}, MLGray, semithick}]
    \node[nninp] (x1) at (0, 1.1) {$x_1$};
    \node[nninp] (x2) at (0, 0) {$x_2$};
    \node[nninp] (x3) at (0,-1.1) {$x_3$};
    \node[nnunit] (n) at (3.4, 0) {$\gsig(z)$};
    \node (a) at (5.6, 0) {$a$};
    \foreach \i in {1,2,3}
      \draw[nnsyn] (x\i) -- node[pos=0.45, above, sloped, font=\scriptsize] {$w_\i$} (n);
    \draw[nnsyn] (n) -- (a);
    \node[font=\scriptsize, MLGray] at (3.4,-1.15) {$z = \vw\cdot\vx + b$};
    \node[font=\scriptsize, MLGray, align=center] at (0,-2.0) {inputs\\(dendrites)};
    \node[font=\scriptsize, MLGray, align=center] at (3.4,-2.0) {sum and squash\\(cell body)};
    \node[font=\scriptsize, MLGray, align=center] at (5.6,-2.0) {activation\\(axon)};
  \end{tikzpicture}
  \caption{An artificial neuron. Inputs are weighted and summed with a bias, and
  the sum is passed through an activation function. The labels in brackets are
  the biological analogy, which is loose: real neurons communicate in timed
  spikes, have thousands of connections with rich internal dynamics, and learn
  by mechanisms that are still not fully understood.}
  \label{fig:neuron}
\end{figure}

\begin{pitfall}{Do not argue from biology}
The analogy was historically productive and is now mostly a source of confusion.
We do not know how the brain learns, and modern networks owe their success to
mathematics, data and computing power, not to fidelity to neuroscience. A design
decision justified by ``that is how the brain does it'' should be treated as
unsupported until it is shown to work on data.
\end{pitfall}

Interest in neural networks rose and collapsed twice --- once after the
perceptron of \cref{ch:ch06} met its limits, once in the 1990s --- before the
third rise, after about 2010, made them the dominant method for speech, images
and text. Two things changed. Digitised records and the web made labelled data
abundant; and graphics processors made the matrix arithmetic of large networks
cheap. Classical models such as logistic regression improve with more data and
then plateau. Large networks keep improving, so for the first time there was a
family of models that could turn ever more data into ever better predictions.

\section{From Features to Learned Features}
\label{sec:learned-features}

Consider a retailer trying to predict whether a new T-shirt will be a top seller.
A single logistic unit on the price alone can do this badly. Suppose instead we
have four inputs --- price, shipping cost, marketing spend and material quality
--- and we reason that sales depend on three intermediate quantities:
\emph{affordability} (from price and shipping), \emph{awareness} (from
marketing) and \emph{perceived quality} (from price and material). A logistic
unit for each intermediate quantity, followed by a logistic unit that combines
them, is a neural network.

\begin{figure}[htbp]
  \centering
  \begin{tikzpicture}[x=3.0cm, y=1.35cm,
      nninp/.style={circle, draw=MLTeal, fill=MLTeal!12, thick, minimum size=6mm},
      nnhid/.style={circle, draw=MLBlue, fill=MLBlue!8, thick, minimum size=7mm},
      nnout/.style={circle, draw=MLAmber, fill=MLAmber!15, thick, minimum size=7mm},
      nnsyn/.style={-{Stealth[length=1.6mm]}, MLGray!70, semithick}]
    \foreach \lab [count=\i] in {price, shipping, marketing, material} {
      \node[nninp] (x\i) at (0,-\i) {};
      \node[left=2pt, font=\scriptsize] at (x\i.west) {\lab};
    }
    \foreach \lab [count=\j] in {affordability, awareness, perceived quality} {
      \node[nnhid] (h\j) at (1,-\j-0.5) {};
      \node[above=1pt, font=\scriptsize, MLBlue, fill=white, inner sep=1pt] at (h\j.north) {\lab};
    }
    \node[nnout] (o) at (2,-2.5) {};
    \node[right=4pt, font=\scriptsize] at (o.east) {$\Prob(\text{top seller})$};
    \foreach \i in {1,...,4} \foreach \j in {1,2,3} \draw[nnsyn] (x\i) -- (h\j);
    \foreach \j in {1,2,3} \draw[nnsyn] (h\j) -- (o);
    \node[font=\scriptsize, MLGray] at (0,-4.9) {input layer $\vx$};
    \node[font=\scriptsize, MLGray] at (1,-5.0) {hidden layer $\va^{[1]}$};
    \node[font=\scriptsize, MLGray] at (2,-5.0) {output layer $a^{[2]}$};
  \end{tikzpicture}
  \caption{A $4\to3\to1$ network for demand prediction. The names on the hidden
  units are what we \emph{hope} they will compute; nothing forces it. Every
  hidden unit sees every input, and training decides which inputs each one
  actually uses.}
  \label{fig:demand-net}
\end{figure}

Two features of \cref{fig:demand-net} are essential. First, the network is
\emph{fully connected}: every hidden unit receives every input. We do not wire
affordability to price by hand; we let training set the weights, and if
marketing spend is irrelevant to affordability its weight will be driven
towards zero. Second, the hidden units' values never appear in the data. The
training set contains only $(\vx, y)$ pairs, which is why the layer is called
\emph{hidden}.

\begin{defbox}{Layers and architecture}
The \term{input layer} is the feature vector $\vx$; it computes nothing. Each
\term{hidden layer} computes a vector of activations from the layer before.
The \term{output layer} produces the prediction. The \term{architecture} is the
set of choices made before training: how many hidden layers, how many units in
each, and which activation function each uses. Training then chooses the weights
and biases. A network of fully connected layers is also called a
\term{multilayer perceptron}.
\end{defbox}

\begin{keyinsight}{Feature engineering versus feature learning}
In \cref{sec:feature-engineering} you decided that the useful feature was
$\text{frontage}\times\text{depth}$, computed it, and handed it to a linear
model. A network's hidden layer computes features too, but what they are is
chosen by gradient descent to make the output layer's job easy. The output layer
is still a logistic regression --- on learned features instead of designed
ones. What remains for the designer is the architecture.
\end{keyinsight}

The same idea scales to inputs nobody can engineer. A greyscale image of
$1000\times1000$ pixels is a vector in $\Real{1{,}000{,}000}$, and no single pixel
says anything about whose face it shows. When researchers inspect networks
trained to recognise faces, they typically find that the first hidden layer
responds to short edges at various orientations, the second to combinations of
edges --- the corner of an eye, the curve of a nostril --- and the third to
coarse whole-face shapes. Train the identical architecture on photographs of
cars, and the intermediate layers learn wheels and windows instead. The
architecture is generic; the features come from the data.

\begin{pitfall}{Interpretations are post hoc}
``Layer two detects eyes'' is a description attached after training, by finding
which inputs excite a unit most strongly. It is not a guarantee, many units have
no tidy description, and a different random initialisation can organise the
same layer quite differently while reaching the same accuracy.
\end{pitfall}

\section{A Layer of Logistic Units}
\label{sec:layer}

A layer is several units that all receive the same input vector. Each has its
own weights and bias, so each computes a different function of the same input.

\subsection{Notation}

A network has several layers, each with several units, applied to many training
examples, so a quantity can need three indices at once. We use the convention of
the companion lecture slides:

\begin{defbox}{Layer notation}
\begin{itemize}
  \item A superscript in square brackets, $[l]$, refers to \textbf{layer} $l$.
  \item A superscript in round brackets, $(i)$, refers to \textbf{training
        example} $i$, as before.
  \item A subscript $j$ refers to \textbf{unit} $j$ within a layer.
\end{itemize}
So $a^{[2](7)}_3$ is the activation of unit $3$ of layer $2$ on example
$7$. By convention the input is layer $0$: $\va^{[0]} = \vx$. A network with
$L-1$ hidden layers and one output layer is said to have $L$ layers; the input
layer is not counted.
\end{defbox}

With this notation, the whole of a network's computation is one equation applied
repeatedly.

\begin{keyinsight}{The one equation}
\[
  a^{[l]}_j \;=\; \gsig\bigl(\vw^{[l]}_j\cdot\va^{[l-1]} + b^{[l]}_j\bigr).
\]
Unit $j$ of layer $l$ takes the \emph{entire} activation vector of the previous
layer, forms a weighted sum, adds its bias, and applies the activation function.
Everything in this chapter is this equation, organised.
\end{keyinsight}

\subsection{The layer in matrix form}

Stack the weight vectors of a layer's $n_l$ units as the columns of a matrix:
\[
  \Wl{l} = \bigl[\,\vw^{[l]}_1\ \ \vw^{[l]}_2\ \ \cdots\ \ \vw^{[l]}_{n_l}\,\bigr]
  \in \Real{n_{l-1}\times n_l},
  \qquad
  \bl{l} = \bigl(b^{[l]}_1, \dots, b^{[l]}_{n_l}\bigr).
\]
Treat an activation vector as a row. Then all $n_l$ weighted sums are one
vector--matrix product, and the activation function is applied entrywise:
\[
  \mathbf{z}^{[l]} = \va^{[l-1]}\Wl{l} + \bl{l}, \qquad
  \va^{[l]} = \gsig\bigl(\mathbf{z}^{[l]}\bigr).
\]
Entry $j$ of $\va^{[l-1]}\Wl{l}$ is $\va^{[l-1]}$ dotted with column $j$ of
$\Wl{l}$, which is $\vw^{[l]}_j\cdot\va^{[l-1]}$, as the one equation requires.

The convention --- units as columns, examples as rows --- is the one TensorFlow
uses, and it is chosen so that a batch of examples can be pushed through a layer
in a single product (\cref{sec:forward-prop}). Some texts use the transpose,
with units as rows of $\mathbf{W}$ and activations as column vectors. The two
are equivalent; mixing them is the most common source of shape errors in a first
implementation.

\subsection{Counting parameters}

Each of a layer's $n_l$ units has one weight per input and one bias, so a fully
connected layer with $n_{\text{in}}$ inputs and $n_{\text{out}}$ units has
\[
  n_{\text{in}}\,n_{\text{out}} + n_{\text{out}}
\]
parameters. The demand network of \cref{fig:demand-net} has $4\cdot3+3 = 15$
in its hidden layer and $3\cdot1+1 = 4$ in its output layer, $19$ in all.

\begin{pitfall}{Count the inputs of \emph{this} layer}
The output layer of the demand network sees the $3$ activations of the hidden
layer, not the $4$ original features. Using the number of features $n$ for every
layer is the most frequent error in parameter counting.
\end{pitfall}

The count grows with the \emph{product} of adjacent widths. A first hidden layer
of $128$ units on a $28\times28$ image needs $784\cdot128+128 = 100{,}480$
parameters before the network has done anything clever, which is part of why
vision networks use the convolutional layers previewed in \cref{ch:ch09}.

\section{Forward Propagation}
\label{sec:forward-prop}

\begin{defbox}{Forward propagation}
\term{Forward propagation} computes a network's output by applying the layer
equation to each layer in turn, from the input to the output:
\[
  \va^{[0]} = \vx, \qquad
  \va^{[l]} = \gsig\bigl(\va^{[l-1]}\Wl{l} + \bl{l}\bigr),\quad l = 1,\dots,L,
  \qquad f(\vx) = \va^{[L]} .
\]
It is all that \term{inference} --- prediction with a trained network --- requires.
\end{defbox}

Stacking layers is composing functions,
$f(\vx) = \gsig^{[L]}\bigl(\cdots\gsig^{[2]}\bigl(\gsig^{[1]}(\vx)\bigr)\cdots\bigr)$,
with each layer's output the next layer's input.

\begin{algorithmbox}{Forward propagation}
\begin{lstlisting}[style=pseudocode]
input: x (row, length n0); W[1..L]; b[1..L]
a <- x
for l = 1 to L:
    # (1 x n_{l-1}) (n_{l-1} x n_l) -> (1 x n_l)
    z <- a @ W[l] + b[l]
    a <- g(z)                   # entrywise
return a                        # f(x)
\end{lstlisting}
Cost: one matrix--vector product per layer, $\sum_l n_{l-1}n_l$ multiplications
in all --- the same as the number of weights.
\end{algorithmbox}

\begin{workedexample}{A $2\to3\to1$ network by hand}
Take the network N231 with sigmoid activations throughout and the hand-chosen
parameters
\[
  \Wl{1} = \begin{bmatrix} 1 & -3 & 5\\ 2 & 4 & -6\end{bmatrix},\quad
  \bl{1} = (-1,\ 1,\ 2),\qquad
  \Wl{2} = \begin{bmatrix} 2\\ -1\\ -3\end{bmatrix},\quad
  b^{[2]} = 0.5 .
\]
Column $j$ of $\Wl{1}$ is the weight vector of hidden unit $j$. Propagate
$\vx = (1, 2)$:

\begin{center}
\small
\renewcommand{\arraystretch}{1.2}
\begin{tabular}{@{}lll@{}}
\toprule
unit & $z = \vw\cdot\va_{\text{in}} + b$ & $a = \gsig(z)$\\
\midrule
$a^{[1]}_1$ & $1\cdot1 + 2\cdot2 - 1 = 4$     & $0.9820$\\
$a^{[1]}_2$ & $-3\cdot1 + 4\cdot2 + 1 = 6$    & $0.9975$\\
$a^{[1]}_3$ & $5\cdot1 - 6\cdot2 + 2 = -5$    & $0.0067$\\
\midrule
$a^{[2]}$   & $2(0.9820) - 0.9975 - 3(0.0067) + 0.5 = 1.4464$ & $\mathbf{0.8094}$\\
\bottomrule
\end{tabular}
\end{center}

$a^{[2]} = 0.8094 \ge 0.5$, so $\yhat = 1$. For $\vx = (-1,-1)$ the same steps
give $\mathbf{z}^{[1]} = (-4,\ 0,\ 3)$,
$\va^{[1]} = (0.0180,\ 0.5000,\ 0.9526)$, $z^{[2]} = -2.8217$ and
$a^{[2]} = 0.0562$, so $\yhat = 0$.

Useful anchors for checking a chain of sigmoids by hand: $\gsig(0) = 0.5$
exactly, $\gsig(4) \approx 0.982$, $\gsig(-5) \approx 0.007$.
\end{workedexample}

\subsection{A whole batch at once}

Put $m$ examples in the rows of a matrix, $\mathbf{A}^{[0]} = \Xmat \in
\Real{m\times n_0}$. The layer equation applies to every row at once:
\[
  \mathbf{Z}^{[l]} = \mathbf{A}^{[l-1]}\Wl{l} + \bl{l},\qquad
  \mathbf{A}^{[l]} = \gsig\bigl(\mathbf{Z}^{[l]}\bigr),
\]
with $\bl{l}$ added to every row (NumPy's broadcasting does this). The shapes
are $(m\times n_{l-1})(n_{l-1}\times n_l)\to(m\times n_l)$: the number of rows
--- examples --- never changes, and the number of columns is the width of the
current layer.

\begin{workedexample}{Both examples in one product}
For N231 with $\Xmat$ holding both examples as rows,
\[
  \Xmat\Wl{1} = \begin{bmatrix}1&2\\-1&-1\end{bmatrix}
                \begin{bmatrix}1&-3&5\\2&4&-6\end{bmatrix}
              = \begin{bmatrix}5&5&-7\\-3&-1&1\end{bmatrix},
  \qquad
  \mathbf{Z}^{[1]} = \Xmat\Wl{1} + \bl{1}
    = \begin{bmatrix}4&6&-5\\-4&0&3\end{bmatrix}.
\]
Row $1$ is the $z$ column of the previous table; row $2$ is the second example.
Then
\[
  \mathbf{A}^{[1]} = \begin{bmatrix}0.9820&0.9975&0.0067\\0.0180&0.5000&0.9526\end{bmatrix},
  \qquad
  \mathbf{A}^{[2]} = \gsig\!\left(\begin{bmatrix}1.4464\\-2.8217\end{bmatrix}\right)
                   = \begin{bmatrix}0.8094\\0.0562\end{bmatrix}.
\]
Shapes: $(2\times2)(2\times3)\to(2\times3)$, then $(2\times3)(3\times1)\to(2\times1)$.
\end{workedexample}

\begin{keyinsight}{Why the matrix form matters}
The loop over units and the matrix product perform identical arithmetic. The
difference is that one product hands the work to compiled linear-algebra code
that uses every processor core and the processor's vector instructions --- and,
on a graphics card, thousands of parallel arithmetic units. This, and not any
new mathematics, is what made networks with millions of parameters practical.
\end{keyinsight}

\subsection{A deeper network}

Nothing limits the number of layers. A handwritten-digit classifier that decides
between $0$ and $1$ from an $8\times8$ image might use $64\to25\to15\to1$:
\[
  \va^{[1]} \in \Real{25},\qquad \va^{[2]} \in \Real{15},\qquad a^{[3]} \in (0,1),
\]
with $64\cdot25+25 = 1625$, $25\cdot15+15 = 390$ and $15\cdot1+1 = 16$
parameters, $2031$ in all. Layers usually narrow towards the output, each
summarising the one before in fewer numbers.

\section{What a Hidden Layer Buys: XOR}
\label{sec:xor}

The clearest demonstration of what a hidden layer adds is a problem that a
single unit provably cannot solve.

\begin{defbox}{Exclusive-or}
The function $\mathrm{XOR}(x_1, x_2)$ on binary inputs is $1$ when exactly one
input is $1$:
\[
  \mathrm{XOR}(0,0) = 0,\quad \mathrm{XOR}(0,1) = 1,\quad
  \mathrm{XOR}(1,0) = 1,\quad \mathrm{XOR}(1,1) = 0 .
\]
\end{defbox}

A single logistic unit has a linear decision boundary, and no line separates
$\{(0,1),(1,0)\}$ from $\{(0,0),(1,1)\}$: the positives lie on one diagonal of
the square and the negatives on the other. Fitting a logistic regression to the
four points confirms it in the most emphatic way possible. The best fit is
$\vw = \mathbf{0}$, $b = 0$, which outputs $0.5$ for every input: by symmetry,
any tilt of the boundary that helps one point hurts another equally, so the
optimum is to commit to nothing.

\begin{workedexample}{A $2\to2\to1$ network for XOR, by hand}
Write XOR as a combination of functions a single unit \emph{can} compute:
\[
  \mathrm{XOR}(x_1,x_2) = \mathrm{AND}\bigl(\mathrm{OR}(x_1,x_2),\
                                           \mathrm{NAND}(x_1,x_2)\bigr).
\]
Each piece is linearly separable, so each is one sigmoid unit with large weights.
Choose
\[
  \text{OR: } \vw = (20, 20),\ b = -10; \qquad
  \text{NAND: } \vw = (-20,-20),\ b = 30; \qquad
  \text{AND: } \vw = (20, 20),\ b = -30 .
\]
The first two form the hidden layer, the third the output unit. Propagate all
four inputs, using $\gsig(10) = 0.999955$ and $\gsig(-10) = 0.000045$:

\begin{center}
\small
\setlength{\tabcolsep}{5pt}
\renewcommand{\arraystretch}{1.2}
\begin{tabular}{@{}ccccccc@{}}
\toprule
$\vx$ & $\mathbf{z}^{[1]}$ & $a^{[1]}_1$ (OR) & $a^{[1]}_2$ (NAND) & $z^{[2]}$ & $a^{[2]}$ & $\yhat$\\
\midrule
$(0,0)$ & $(-10,\ 30)$ & $0.000045$ & $1.000000$ & $-9.9991$ & $0.000045$ & $0$\\
$(0,1)$ & $(10,\ 10)$  & $0.999955$ & $0.999955$ & $9.9982$  & $0.999955$ & $1$\\
$(1,0)$ & $(10,\ 10)$  & $0.999955$ & $0.999955$ & $9.9982$  & $0.999955$ & $1$\\
$(1,1)$ & $(30,\ -10)$ & $1.000000$ & $0.000045$ & $-9.9991$ & $0.000045$ & $0$\\
\bottomrule
\end{tabular}
\end{center}

All four correct, each with a margin of about $10$ on the logit scale.
\end{workedexample}

\begin{figure}[htbp]
  \centering
  \begin{minipage}[t]{0.48\linewidth}
    \centering
    \begin{tikzpicture}
      \begin{axis}[mlaxis, width=4.6cm, height=4.6cm, scale only axis,
        xlabel={$x_1$}, ylabel={$x_2$}, xmin=-0.3, xmax=1.3, ymin=-0.3, ymax=1.3,
        title={input space}]
        \addplot[MLBlue, thick, dashed, domain=-0.3:1.3, forget plot] {0.5 - x};
        \addplot[MLAmber, thick, dashed, domain=-0.3:1.3, forget plot] {1.5 - x};
        \addplot[only marks, mark=*, mark size=2.4pt, MLRose, forget plot]
          coordinates {(0,0)(1,1)};
        \addplot[only marks, mark=square*, mark size=2.4pt, MLTeal, forget plot]
          coordinates {(0,1)(1,0)};
        \node[font=\tiny, MLBlue, rotate=-45] at (axis cs:0.05,0.62) {OR};
        \node[font=\tiny, MLAmber, rotate=-45] at (axis cs:0.55,1.12) {NAND};
      \end{axis}
    \end{tikzpicture}
  \end{minipage}\hfill
  \begin{minipage}[t]{0.48\linewidth}
    \centering
    \begin{tikzpicture}
      \begin{axis}[mlaxis, width=4.6cm, height=4.6cm, scale only axis,
        xlabel={$a^{[1]}_1$ (OR)}, ylabel={$a^{[1]}_2$ (NAND)},
        xmin=-0.3, xmax=1.3, ymin=-0.3, ymax=1.3, title={hidden space}]
        \addplot[MLGray, thick, domain=-0.3:1.3, forget plot] {1.5 - x};
        \addplot[only marks, mark=*, mark size=2.4pt, MLRose, forget plot]
          coordinates {(0.000045,1.0)(1.0,0.000045)};
        \addplot[only marks, mark=square*, mark size=2.4pt, MLTeal, forget plot]
          coordinates {(0.999955,0.999955)};
        \node[font=\tiny, MLGray, anchor=west] at (axis cs:0.02,0.3) {AND: $a_1 + a_2 = 1.5$};
      \end{axis}
    \end{tikzpicture}
  \end{minipage}
  \caption{What the hidden layer does to XOR. Left: in input space the
  positives (squares) and negatives (circles) sit on opposite diagonals; no line
  separates them. The two hidden units draw two lines, one each. Right: the
  same four inputs mapped to the hidden layer's activations. The two positives
  land on the same point, near $(1,1)$; the negatives land near $(0,1)$ and
  $(1,0)$. Now a single line --- the output unit --- separates them.}
  \label{fig:xor}
\end{figure}

\begin{keyinsight}{A hidden layer re-represents the input}
The output unit is still a linear classifier. What changed is its input: the
hidden layer mapped the four points to a space in which the classes are linearly
separable. That is the general job of hidden layers --- to transform the data
until the last layer's simple linear decision is enough. Here we chose the
transformation by hand; \cref{ch:ch09} shows how gradient descent finds one from
data.
\end{keyinsight}

\begin{researchspotlight}{McCulloch \& Pitts (1943) --- A logical calculus of nervous activity}
\textbf{The problem.} Can networks of simple, idealised neurons compute
anything interesting, and if so, what?

\textbf{The idea.} Model a neuron as a threshold unit: it fires if the weighted
sum of its binary inputs exceeds a threshold. Show that single units implement
AND, OR and NOT, and that networks of them implement any finite logical
expression.

\textbf{The result.} Networks of threshold units are as expressive as Boolean
logic --- the first formal model of neural computation, and an influence on the
design of the stored-program computer. The paper contains no learning at all:
every weight is set by hand.

\textbf{Why it is here.} The XOR network of \cref{sec:xor} is a McCulloch--Pitts
construction with the hard threshold softened into a sigmoid. The softening is
everything: a threshold has zero derivative almost everywhere, so its weights
cannot be learned by gradient descent, while the sigmoid's can.
\parencite{mcculloch1943}
\end{researchspotlight}

\section{Why Nonlinear Activations Are Required}
\label{sec:collapse}

The activation function has been the sigmoid throughout. It is natural to ask
whether it is needed at all, since the weighted sums already do most of the
arithmetic. It is, and the reason takes two lines.

\begin{derivation}{Linear layers compose to one linear layer}
Suppose every layer uses the identity activation $\gsig(z) = z$. Two layers give
\[
  \va^{[2]} = \bigl(\vx\Wl{1} + \bl{1}\bigr)\Wl{2} + \bl{2}
            = \vx\,\underbrace{\Wl{1}\Wl{2}}_{\mathbf{W}'}
              + \underbrace{\bl{1}\Wl{2} + \bl{2}}_{\mathbf{b}'} ,
\]
which is a single linear layer with parameters $(\mathbf{W}',\mathbf{b}')$. By
induction on the number of layers, any depth collapses the same way. A network
of linear layers with a linear output is linear regression; with a sigmoid
output it is logistic regression. Depth adds parameters but no ability.
\end{derivation}

\begin{workedexample}{N231 with every activation removed}
Replace every sigmoid in N231 by the identity. Then
\[
  \mathbf{W}' = \Wl{1}\Wl{2}
    = \begin{bmatrix}1\cdot2 - 3\cdot(-1) + 5\cdot(-3)\\ 2\cdot2 + 4\cdot(-1) - 6\cdot(-3)\end{bmatrix}
    = \begin{bmatrix}-10\\ 18\end{bmatrix},
  \qquad
  b' = (-1)(2) + (1)(-1) + (2)(-3) + 0.5 = -8.5 .
\]
For $\vx = (1,2)$ the two-layer computation gives $17.5$ and
$-10 + 36 - 8.5 = 17.5$; for $(-1,-1)$, both give $-16.5$. Thirteen parameters
reproduce a function that three would describe.
\end{workedexample}

The non-linearity is what lets a network represent functions a linear model
cannot. How much it can represent is the subject of a remarkable theorem.

\begin{researchspotlight}{Cybenko (1989) --- Universal approximation}
\textbf{The problem.} Is there any continuous function that a network with one
hidden layer of sigmoid units cannot represent, however many units it has?

\textbf{The idea.} Treat the set of functions
$\sum_{j} v_j\,\gsig(\vw_j\cdot\vx + b_j)$ as a subset of the continuous
functions on a bounded region, and use functional analysis --- the Hahn--Banach
and Riesz representation theorems --- to show that nothing can be bounded away
from it.

\textbf{The result.} Finite sums of sigmoids are dense in the continuous
functions on the unit cube: any continuous function can be approximated to any
accuracy by a network with one hidden layer, provided the layer is wide enough.
Later work extended the result to essentially any non-polynomial activation.

\textbf{Why it is here.} It says a network \emph{can} represent the function
you want. It does not say how many units that takes --- often exponentially
many in the input dimension --- nor that gradient descent will find the weights,
nor that the result will generalise from finite data. Those three gaps are why
depth, training and regularization still matter, and why the theorem is a
licence to try networks rather than a guarantee that they will work.
\parencite{cybenko1989}
\end{researchspotlight}

\section{Networks in Code}
\label{sec:nn-code}

A practitioner writes very little code to build a network. The cost of that
brevity is that, when something goes wrong, only someone who knows what the
code computes can diagnose it. We write N231 three ways.

\subsection{In TensorFlow}

\begin{lstlisting}[style=tfkeras, caption={N231 in Keras, with the hand-chosen weights loaded.}, label={lst:n231-keras}]
import numpy as np
from tensorflow.keras import Sequential, Input
from tensorflow.keras.layers import Dense

model = Sequential([
    Input(shape=(2,)),                       # two features
    Dense(units=3, activation='sigmoid'),    # hidden layer
    Dense(units=1, activation='sigmoid'),    # output layer
])
W1 = np.array([[1., -3., 5.], [2., 4., -6.]])
b1 = np.array([-1., 1., 2.])
W2 = np.array([[2.], [-1.], [-3.]])
b2 = np.array([0.5])
model.layers[0].set_weights([W1, b1])
model.layers[1].set_weights([W2, b2])

X = np.array([[1., 2.], [-1., -1.]])   # one row per example
print(model.predict(X))                # [[0.8094] [0.0562]]
\end{lstlisting}

A \code{Dense} layer is a fully connected layer: \code{units} sets its width and
\code{activation} its $\gsig$. Its weights are stored exactly as in
\cref{sec:layer}, as a matrix of shape $(n_{\text{in}}, \text{units})$ whose
columns are the units, which is why \code{set\_weights} accepts \code{W1} as
written. Run against the NumPy computation of \cref{sec:forward-prop}, the two
agree to $3\times10^{-9}$ --- the difference being TensorFlow's single-precision
arithmetic.

\subsection{In NumPy, from scratch}

\begin{lstlisting}[style=pythonml, caption={Forward propagation from scratch: a loop over units, then a loop over layers.}, label={lst:dense-loop}]
def g(z):
    return 1 / (1 + np.exp(-z))

def dense(a_in, W, b):
    """a_in (n_in,), W (n_in, units), b (units,)"""
    units = W.shape[1]
    a_out = np.zeros(units)
    for j in range(units):
        w = W[:, j]          # column j = unit j's weights
        a_out[j] = g(np.dot(w, a_in) + b[j])
    return a_out

def sequential(x, params):
    a = x
    for W, b in params:      # one layer per (W, b) pair
        a = dense(a, W, b)
    return a

params = [(W1, b1), (W2, b2)]
sequential(np.array([1., 2.]), params)   # array([0.8094])
\end{lstlisting}

\subsection{In NumPy, vectorised}

\begin{lstlisting}[style=pythonml, caption={The vectorised layer: one matrix product for all units and all examples.}, label={lst:dense-vec}]
def dense_v(A_in, W, b):
    """A_in (m, n_in), W (n_in, units) -> (m, units)"""
    return g(A_in @ W + b)     # b is added to every row

A1 = dense_v(X, W1, b1)               # (2, 3)
A2 = dense_v(A1, W2, b2)    # (2, 1): [[0.8094], [0.0562]]
\end{lstlisting}

\begin{pitfall}{One example is a $1\times n$ matrix, not a vector}
\code{np.array([200., 17.])} has shape $(2,)$ --- a 1-D array with neither rows
nor columns. \code{np.array([[200., 17.]])} has shape $(1,2)$: one example with
two features. \code{np.array([[200.], [17.]])} has shape $(2,1)$: two examples
with one feature each. TensorFlow layers expect the second form; passing the
first raises an error or, worse, is broadcast into something unintended, and the
third silently means a different dataset. A bias of shape $(\text{units},1)$
instead of $(\text{units},)$ has the same effect inside \code{dense\_v}. Print
\code{.shape} whenever in doubt, and \code{assert} it in any function that takes
a batch.
\end{pitfall}

\begin{ethicsnote}{Capability is not understanding}
A network that classifies chest X-rays as well as a specialist on a benchmark
has learned a mapping from pixels to labels on that benchmark's distribution. It
has not learned anatomy, and it may be relying on cues --- a scanner's
watermark, the hospital that produced the image --- that a specialist would never
use and that will not be present elsewhere. Claims about what a network
``knows'' or ``understands'' should be translated into testable statements:
\emph{this} model achieves \emph{this} error on \emph{this} population. The
evaluation methods of \cref{ch:ch10} and the subgroup checks of \cref{ch:ch11}
are how such statements are tested.
\end{ethicsnote}

\begin{labbox}{Lab 8 --- Neurons, layers and forward propagation}
Build single neurons in TensorFlow with linear and sigmoid activations, load
them with the parameters fitted in earlier chapters, and confirm that they
reproduce those chapters' predictions. Build N231, load the weights of
\cref{sec:forward-prop}, and reproduce $0.8094$ and $0.0562$ in Keras, in the
looped NumPy implementation, and in the vectorised one.

Then build the XOR network of \cref{sec:xor} and reproduce its table. Fit a
single logistic unit to the XOR points and confirm that it learns $\vw =
\mathbf{0}$. Finally, count the parameters of $64\to25\to15\to1$ by formula and
with \code{model.summary()}: both must give $2031$.
\end{labbox}

\section*{Chapter Summary}
\addcontentsline{toc}{section}{Chapter Summary}

\begin{itemize}
  \item A unit computes $a = \gsig(\vw\cdot\va + b)$; with the sigmoid it is a
        logistic regression. The biological analogy motivated the model and
        should not be used to justify design decisions.
  \item A hidden layer computes features that training, not the designer,
        chooses. The output layer is a logistic regression on those features.
        The designer chooses the architecture.
  \item Notation: $[l]$ is the layer, $(i)$ the example, $j$ the unit;
        $\va^{[0]} = \vx$. Every layer computes
        $a^{[l]}_j = \gsig(\vw^{[l]}_j\cdot\va^{[l-1]} + b^{[l]}_j)$.
  \item Stacking unit weights as columns gives $\Wl{l}$ of shape
        $(n_{l-1}, n_l)$ and the layer $\va^{[l]} = \gsig(\va^{[l-1]}\Wl{l} +
        \bl{l})$. A layer has $n_{\text{in}}n_{\text{out}} + n_{\text{out}}$
        parameters.
  \item Forward propagation applies the layer equation from input to output. On
        a batch, $\mathbf{A}^{[l]} = \gsig(\mathbf{A}^{[l-1]}\Wl{l} + \bl{l})$.
        N231 gives $0.8094$ for $(1,2)$ and $0.0562$ for $(-1,-1)$.
  \item A single unit cannot compute XOR --- the best logistic fit outputs $0.5$
        everywhere --- but two layers can, because the hidden layer maps the
        inputs to a space where they are linearly separable.
  \item Without a non-linear activation, any depth collapses to one linear
        layer: N231 becomes $-10x_1 + 18x_2 - 8.5$.
  \item One wide hidden layer can approximate any continuous function, which
        says nothing about how wide, how to find the weights, or whether the
        result generalises.
  \item Keras, a loop, and a matrix product compute the same function, provided
        examples are rows and units are columns throughout.
\end{itemize}

\section*{Exercises}
\addcontentsline{toc}{section}{Exercises}

\exercisesection{Short Questions}

\sq Describe an artificial neuron: its inputs, parameters, computation and
output. Why is a single sigmoid unit the same as logistic regression?

\sq Give two reasons neural networks became dominant after about 2010.

\sq Why is a hidden layer called ``hidden''?

\sq Explain the meaning of each index in $a^{[2](7)}_3$.

\sq What does the designer of a neural network still choose, and what does
training choose?

\sq How many layers does a network with input $\vx$, three hidden layers and one
output unit have, under this chapter's convention?

\sq A layer receives $\mathbf{A}^{[l-1]}$ of shape $(100, 64)$ and has $25$
units. Give the shapes of $\Wl{l}$, $\bl{l}$, $\mathbf{Z}^{[l]}$ and
$\mathbf{A}^{[l]}$.

\sq Why does no single logistic unit compute XOR?

\sq What does the universal approximation theorem guarantee, and what three
things does it not?

\sq State the shapes of \code{np.array([5, 3])}, \code{np.array([[5, 3]])} and
\code{np.array([[5], [3]])}, and say which a Keras layer expects for one example
with two features.

\exercisesection{Analytical Problems}

\ap Show that the entry in row $i$, column $j$ of $\mathbf{A}^{[l-1]}\Wl{l}$ is
$\vw^{[l]}_j\cdot\va^{[l-1](i)}$, and hence that the batch layer equation is the
one equation applied to every example and every unit.

\ap Compute the number of parameters in each layer and in total for networks of
sizes $64\to25\to15\to1$ and $784\to128\to64\to10$. Which layer holds most of
the parameters in each, and why?

\ap Propagate $\vx = (0, 0)$ and $\vx = (0.5, -0.5)$ through N231 by hand,
giving $\mathbf{z}^{[1]}$, $\va^{[1]}$, $z^{[2]}$ and $a^{[2]}$ for each, and
state the predictions.

\ap Using the matrix form, compute $\mathbf{Z}^{[1]}$ and $\mathbf{A}^{[2]}$ for
N231 with $\Xmat = \begin{bmatrix}0&0\\2&1\end{bmatrix}$.

\ap Prove by induction that a network of $L$ layers with identity activations
computes an affine function of $\vx$, and give $\mathbf{W}'$ and $\mathbf{b}'$
in terms of the layers' parameters.

\ap Prove that no choice of $(w_1, w_2, b)$ gives a logistic unit that outputs
at least $0.5$ on $(0,1)$ and $(1,0)$ and less than $0.5$ on $(0,0)$ and $(1,1)$.
(Hint: add the two pairs of inequalities on $z$.)

\ap Design by hand a $2\to2\to1$ sigmoid network computing XNOR, the negation of
XOR, and verify it on all four inputs.

\exercisesection{Design Problems}

\dpr A microcontroller can store at most $5{,}000$ parameters. Design a network
for $28\times28$ greyscale images that decides between two classes and fits the
budget. Show the count, and say what you give up to fit.

\dpr You inherit a function that claims to perform forward propagation. Design
three hand-computable test cases that would expose, respectively, a missing bias,
the wrong input to the second layer, and a missing activation function. State
the expected output of each.

\dpr A colleague claims hidden unit $7$ of a trained X-ray classifier ``detects
fractures''. Design an experiment to support or refute the claim, and state
what result would refute it.

\dpr A mobile app must score two million transactions overnight with a
$30\to16\to8\to1$ network. Design the inference code: the batch size, the
memory needed per batch, and the shape of every intermediate matrix. Justify the
batch size.

\exercisesection{Programming Exercises}

\pe Implement \code{dense}, \code{sequential} and \code{dense\_v} as in
\cref{sec:nn-code}, and verify that all three reproduce the N231 outputs of
\cref{sec:forward-prop}.

\pe Build N231 in Keras, load the weights, and write an assertion-based test
that your NumPy implementation agrees with \code{model.predict} on $1000$ random
inputs to within $10^{-6}$. Explain the tolerance.

\pe Time the looped and vectorised layers on a $(10{,}000, 784)$ input and a
$784\to128$ layer. Report the speed-up on your machine, and explain why it is so
large although the arithmetic is the same.

\pe Train a single logistic unit and a $2\to2\to1$ network on the four XOR points
with Keras. Report whether each learns XOR, and, for the network, plot the four
points in hidden space after training. Repeat with five random seeds; does
training always succeed?

\exercisesection{Self-Test}

\qz A single sigmoid unit is equivalent to: (a)~linear regression; (b)~logistic
regression; (c)~a decision tree; (d)~$k$-means.

\qz In $a^{[l]}_j$, the superscript $[l]$ denotes: (a)~the training example;
(b)~the layer; (c)~the unit; (d)~the iteration.

\qz A layer with $64$ inputs and $25$ units has: (a)~$64$ parameters; (b)~$1600$;
(c)~$1625$; (d)~$89$.

\qz In the convention of this chapter, the weight matrix of a layer has shape:
(a)~$(\text{units}, n_{\text{in}})$; (b)~$(n_{\text{in}}, \text{units})$;
(c)~$(m, \text{units})$; (d)~$(\text{units}, \text{units})$.

\qz Forward propagation is required for: (a)~inference only; (b)~training only;
(c)~both inference and training; (d)~neither.

\qz A network whose layers all use the identity activation is equivalent to:
(a)~a deeper network; (b)~a single linear layer; (c)~a decision tree; (d)~XOR.

\qz The best single logistic unit fitted to XOR outputs: (a)~$0$ everywhere;
(b)~$1$ everywhere; (c)~$0.5$ everywhere; (d)~XOR.

\qz The universal approximation theorem guarantees that: (a)~gradient descent
finds the best weights; (b)~one sufficiently wide hidden layer can approximate
any continuous function; (c)~networks always generalise; (d)~deeper is better.

\furtherreading

\textcite{james2021} §10.1 introduces single-layer networks with exactly this
chapter's reading of hidden units as learned features, and §10.2 the multilayer
network, written with weight matrices and applied to handwritten digits; both
are short and worth reading in full. \textcite{bishop2006} §5.1 gives the
feed-forward network function in general form, including how the bias can be
absorbed into the weights by an extra constant input, a careful statement of the
approximation properties behind the Cybenko spotlight, and a discussion of the
symmetries --- permuting hidden units, flipping signs --- that make many different
weight settings compute the same function.

\bigskip
\noindent\emph{Next:} \Cref{ch:ch09} turns the forward pass into a learning
algorithm: the loss, the activation functions beyond the sigmoid, softmax for
many classes, and back-propagation, which computes every gradient of a network in
one backward sweep.
